\subsection{Two minimal rows, one extra actor, and at most five states}

\begin{lemma}[Proved: pure two-row bands on at most five states]
\label{CP2:A-pure-two-row-five-states-products}
Let $P_c$ be the identity-adjoined $2\times c$ rectangular
band, with $c\ge2$.  Every finite left-$R_3$, right-$P_c$
biset of at most five states has zero rational degree-two
products on its collage.  No Betti-number hypothesis is needed.
\end{lemma}
\begin{proof}
The empty biset is immediate from the pure-band projective
resolution.  Otherwise put $T=SK$, the union of the minimal
band images.  As a right $P_c$ set, $T$ is a disjoint union
of column-quotient orbits $O_\pi$, indexed by partitions
$\pi$ of the $c$ column labels.  A map of two such orbits
exists exactly when the target partition is coarser, and
is then unique: it takes the state represented by column
$j$ to the target state represented by the same column.
Every orbit endomorphism is therefore the identity.

Let $W$ be the augmentation kernel of the rational column
permutation module, and let
\[
 W_\pi=\ker(W\longrightarrow
            \ker(\mathbb QO_\pi\to\mathbb Q)).
\]
Thus $W_\pi$ is the direct sum of the augmentation-zero
spaces on the individual blocks of $\pi$.
The complete length-two pure-band resolution gives naturally
\[
 \operatorname{Tor}^{\mathbb QP_c}_2(\mathbb Q[S],\mathbb Q)
        =\bigoplus_{O_\pi\subseteq T} W_\pi.                 \tag{1}
\]
Indeed its top differential is the literal column-incidence
map, with one block for each orbit; every state in its
image lies in $T$.  The inclusion of the $T$-complex into
the $S$-complex is injective in degree one and identical
in degree two, so it has exactly the same cycle kernel.
For the core set $T$ the column differences span the entire
degree-one term, giving
\[
 \operatorname{Tor}^{\mathbb QP_c}_1(\mathbb Q[T],\mathbb Q)=0.
                                                               \tag{2}
\]
The right-$R_3$ resolution has length one.  Therefore (1)
and the two-block sequence identify fourth collage homology
with
$\operatorname{Tor}^{\mathbb QR_3}_1
 (\mathbb Q,\bigoplus W_\pi)$.
In particular the actual core inclusion induces an
isomorphism on $H_4$.  It suffices to kill the products
on the core collage.

A one-state core orbit gives an actual $P_c$-fixed arrow,
which kills all products by the fixed-arrow relative
suspension argument.  Assume there is no such orbit.
Every orbit then has at least two states, so there are
at most two core orbits.
The two $R_3$ retractions have common image on their
orbit set.  If there is only one orbit, or both orbits
are fixed, their actions on every orbit are the identity.
The module in (1) is then trivial under $R_3$, so its
first derived coinvariants are zero and $H_4=0$.

The remaining case has one fixed orbit $O_{\pi_0}$
and one orbit $O_{\pi_1}$ mapping to it under both
retractions.  Uniqueness of the orbit map implies that
both retractions are the same map, and $\pi_0$ is
coarser than $\pi_1$.
On (1) the map is the identity on $W_{\pi_0}$ and the
natural inclusion $W_{\pi_1}\hookrightarrow W_{\pi_0}$
on the other summand.  The literal right-zero resolution
therefore gives
\[
 H_4(B\mathcal C_T;\mathbb Q)\cong W_{\pi_1}.                \tag{3}
\]

Let $G$ be the product of the symmetric groups on the
individual blocks of $\pi_1$.  Permuting columns by $G$
is an actual automorphism of the pure-band monoid and
fixes every state of both core orbits, since $\pi_0$ is
coarser.  Thus it acts on the actual core collage.
Its action on (3) is the standard action on $W_{\pi_1}$,
whose invariant subspace is zero: a fixed vector is
constant on each block and also has sum zero on that block.

On the other hand, $G$ acts trivially on $H_2$ of this
collage.  To see this with the actual maps, (2) makes
first bridge homology the first derived $R_3$-coinvariants
of the core-orbit permutation module; $G$ fixes that
module pointwise.  Second bridge homology is the ordinary
coinvariant module of (1), namely $W_{\pi_0}$, whose
map to $H_2(BP_c)=W$ is its literal inclusion.  The
two-block sequence consequently expresses $H_2$ as an
extension of $W/W_{\pi_0}$ by that first bridge group.
Both are trivial $G$ modules.  Averaging a linear section
over the finite group splits this extension equivariantly
over $\mathbb Q$, so $H_2$ is itself a trivial $G$ module.
The same holds dually for $H^2$; the dual of (3) has no
invariants.  Every product of second classes is invariant,
and hence zero.  The actual core inclusion detects $H^4$,
so the conclusion also holds for the original $S$.
\end{proof}

\begin{theorem}[Proved: one additional actor over a two-row band]
\label{CP2:A-two-row-one-extra-five-states}
Let $M$ have proper minimal band $K=\{e_{ij}:i=0,1,
\ 0\le j<c\}$, with $c\ge2$, and at most two elements
outside $K$, including its identity.
For every finite left-$R_3$, right-$M$ biset of at most
five states, all rational degree-two products of the
collage vanish.  No boundedness or Betti-number assumption
is imposed.
\end{theorem}
\begin{proof}
The pure-band case is the preceding lemma.  Otherwise
write $M=\{1,a\}\sqcup K$, and let $\sigma$ and $\tau$
be the row and column actors of $a$.  Its row-augmentation
scalar $\chi$ is $0,1$, or $-1$.

If $\chi=0$, the row actor is constant, say at row $i_*$.
Collapse the image of $\tau$ to one column and keep every
other column distinct.  This defines a column quotient
$\pi:J\to J'$ and a monoid quotient to the pure two-row
band with $|J'|$ columns, sending $a$ to the band element
in row $i_*$ and the collapsed column.
Indeed $a^2$ is either $a$ or a band element, since its
row actor is not invertible.  In the latter case the
column of $a^2$ lies in $\operatorname{im}\tau$; in either
case its image equals the square of the displayed band
element.  The mixed products agree by the definition of
$\pi$, since every $\tau(j)$ lies in the collapsed image.

Write $W=\ker(\mathbb QJ\to\mathbb Q)$.
The balanced corner presentation is $W/W\tau$.
The space $W\tau$ is exactly the augmentation-zero
space on $\operatorname{im}\tau$, which is precisely
the kernel of the actual map to
$\ker(\mathbb QJ'\to\mathbb Q)$.
Thus this corner quotient is an isomorphism on $H_2$.

If $\chi=1$, the row actor is the identity and
$a^2\in\{1,a\}$: a band square would have constant
row action.  Form $J'$ by the equivalence relation
generated by $j\sim\tau(j)$, and send $a$ to the
global identity of the pure-band target.  This is
again an actual monoid quotient.  The kernel on
column augmentations is exactly $W(\tau-1)$,
so the balanced corner $H_2$ map is an isomorphism.

Both quotient constructions preserve the two rows of
every fixed minimal left ideal $Me_{0j}$ bijectively.
Lemma~\ref{CP2:A-idempotent-quotient-bridge-connectivity}
therefore gives bridge $H_1$ isomorphism and $H_2$
surjection.  Together with the directly checked corner
$H_2$ isomorphism, the two-block exact sequence gives
an actual collage $H_2$ isomorphism.
The induced set $S\otimes_M H$ is a quotient of $S$,
so still has at most five states.  Its pure-band collage
has zero products by the preceding lemma; these pull
back to every second class of the original collage.
If $J'$ has one element, its corner $H_2$ is zero,
and the relative suspension argument supplies the
same conclusion directly.

It remains to handle $\chi=-1$.  Then $\sigma$ is the
row transposition.  Its square is the identity, so
$a^2=1$; neither $a$ nor a band element has the
required square row action.  Hence $\tau$ is an
involution.  Choose its transposed pairs $(u_l,v_l)$,
$1\le l\le t$, and put
\[
 R=\mathbb QM,\quad e=e_{00},\quad q=(1-a)/2,
 \quad w=e_{10}-e_{00},\quad
 \delta_l=e_{0u_l}-e_{0v_l}.
\]
There is a complete left-projective resolution
\[
 0\longrightarrow(Re)^t\xrightarrow{d_2}Rq
 \xrightarrow{d_1}Re\longrightarrow\mathbb Q
 \longrightarrow0.                                         \tag{4}
\]
Here $d_1$ is right multiplication by $w$ and the $l$th
part of $d_2$ is right multiplication by $\delta_l$.
To verify it, $qw=w$, and $Rq$ has basis $q$ together
with the $2t$ differences
$e_{iu_l}-e_{iv_l}$ for $i=0,1$.
The map $d_1$ sends $q$ to $w$ and kills those
differences; $d_2$ maps the $2t$ displayed basis
vectors bijectively to them.  This also covers $t=0$.

Let $B=\mathbb QP_c\subset R$ and use its pure-band
resolution, with degree-one term $B(1-e)$ and top
term $(Be)^{c-1}$.  The inclusion comparison is the
identity in degree zero and right multiplication
by $q$ in degree one.  In degree two it is induced
by the antisymmetrizing projection
\[
 W\longrightarrow W^-,\qquad z\longmapsto zq=(z-za)/2,
                                                               \tag{5}
\]
where $W^-$ has basis $\delta_1,\ldots,\delta_t$.
This is a chain map because the pure-band top
differentials are the differences $e_{0j}-e_{00}$.
The comparison does not require $eq=0$; its degree-one
value on $1-e$ is $(1-e)q$.

For $L=\mathbb Q[S]$, a target top cycle is a tuple
$(s_l)\in(Le)^t$ satisfying
$\sum_l s_l\delta_l=0$.
Express each $\delta_l$ in the column-difference
basis $e_{0j}-e_{00}$, $j\ne0$, and use those same
constant coefficients on the tuple $(s_l)$.
The resulting tuple is a cycle in the pure-band
complex, because its boundary is the same displayed
sum.  Equation (5) sends it back to $(s_l)$ since
$\delta_l q=\delta_l$.
This is a natural left-$R_3$ linear section of the
map on second source Tor groups.

Exactly as in the source-$500$ comparison, the
length-one right-$R_3$ resolution makes this a
surjection on third bridge homology.  Both source
corners have projective dimension at most two by
their complete displayed resolutions.  The actual
pure-band inclusion therefore induces a surjection
on fourth collage homology and an injection on
fourth cohomology.  Its pure-band source has at most
five states and zero products by the preceding lemma.
Restriction detects every product, proving the claim.
\end{proof}

\begin{corollary}[Proved: the two-by-four order-ten source family]
\label{CP2:A-ten-two-by-four-five-state-family}
Every eighteen-arrow collage with ten-arrow source corner,
three-arrow opposite corner $R_3$, five connecting arrows,
and a $2\times4$ rectangular minimal source band has zero
rational degree-two products.  In particular no such
collage is a rational $\mathbb CP^2$ model.
\end{corollary}
\begin{proof}
The source has precisely its identity and one other
element outside the eight-element band.  Apply the theorem.
\end{proof}
